\begin{abstract}
Presence in virtual reality is usually measured with a post-session questionnaire, leaving the session itself unmeasured. We ask whether concurrent think-aloud speech, scored by an evidence-gated LLM pipeline, converges with that questionnaire. In a between-subjects study, 32 participants played three VR games, 16 thinking aloud and 16 silently; both groups completed the Igroup Presence Questionnaire after each game. In the think-aloud group, speech-derived scores correlated with the questionnaire on nine held-out participants ($r=.61$, $p<.001$) and, within each participant, identified the game rated most present in 69\% of cases (chance 33\%). Speech scored 0.22 points below the questionnaire, inside its measurement error under an equivalence test. Between groups, the silent group did not report more presence than the think-aloud group; no difference survived multiplicity correction. The method tracks questionnaire presence at the group level and keeps each score traceable to its spoken moment; involvement items, rarely verbalized, remain its weakest point.
\end{abstract} 
% 재훈님, introduction 좀 바꿨는데, 혹시 통화 가능하신가요?  보시고 가능하시면 전화 한번 주세요 :) 내용 맞춰보도록 하겠습니다 

% between subject test 로 진행한 내용이 들어가야할 것 같으며, 
% between subject 에서 not think aloud 그룹에서 모아온 결과가 어떤 것인지에 대한 내용이 포함되어야 할 것 같아요
% 그리고 뒤에 highest presence game 69% 는 무엇을 이야기 하는 내용인지 잘 모르겠습니다. 
% TOST p  에서  TOST 는 무엇인지요..? 
% Think aloud 를 통해서 자동 측정한 결과가      think aloud 하고 설문지 한 결과하고 얼마나 상관관계가 있었는지에 대한 결과와 
% think aloud 결과랑 , think aloud 하지 않고 설문지만 한 그룹의 관계에 대한 결과 비교가 위에 포함되어야 할 것 같습니다. 
% ---- 초록 주석 답변 (2026-09-10, 저자) ----
% [between-subject 설계] 셋째 문장에 "In a between-subjects study … 16 thinking aloud and 16 silently; both groups completed the IPQ after each game" 으로 설계와 두 군 공통 측정을 명시함.
% [not think-aloud 군의 결과] 여섯째 문장 "Between groups, the silent group did not report more presence than the think-aloud group; no difference survived multiplicity correction" — 실측: 발화 군이 IPQ 총점 +0.44 높았으나 (p=.030) Holm 보정 후 7개 지표 모두 비유의 (5.6절, Table tab:group).
% [69% 의 뜻] "identified the game rated most present in 69% of cases (chance 33%)" — 참가자 본인이 설문에서 가장 높게 매긴 게임을 발화 점수가 맞춘 비율 11/16, 우연 기준 33% (5.1절). 76% (게임 쌍 방향 일치) 는 단어 제한으로 초록에서 빼고 본문에만 둠.
% [TOST] 약어 대신 "under an equivalence test" 로 풀어씀 — TOST = two one-sided tests, 동등성 검정 (5.2절). p 값은 본문에.
% [TA 자동 점수 대 TA 군 설문의 상관] 넷째 문장 r=.61, p<.001 (홀드아웃 9명, 5.2절). "Speech"/"Survey" 따옴표 호칭은 정의가 없어 "speech-derived scores"/"the questionnaire" 로 바꿈.
% [TA 군 대 비TA 군 비교] 위 [not think-aloud 군의 결과] 와 같은 문장으로 답함 (5.6절 RQ3).
% 초록 150 단어 (PCS 제한 150). 초록 안의 \cite 는 본문에 근거 절이 있으므로 제거함.